\section{The compact positive-ratio passage}
\label{sec:compact-positive}

We next treat the coefficient face \(\alpha _0=-s^2<0\) with
\(p=\alpha _1/s\) in a fixed interval \([p_0,P]\subset(0,\infty)\).
The low-height orbits pass between a positive equatorial saddle and
a negative saddle escaping to \(X=-\infty\), whose characteristic
number tends to zero.

The escaping saddle prevents a comparison on fixed compact source
sections alone. We parametrize its passage by the point where the strong
leaf meets the equator. An axis integral measures the normal expansion
all the way to that saddle. Matching it with the positive Dulac map gives
a curvature estimate in the input of the same return, uniformly as the
negative characteristic number tends to zero. We then join this endpoint
estimate to the compact middle and the high interval, where the derivative
is bounded.

\begin{lemma}[Domains of the compact positive-ratio passage]
\label{ann:lem-compact-positive-domains}
For \(p_0\le p\le P\), sufficiently small \(s>0,|\gamma|\), and one
fixed sufficiently small bound on \(z(\gamma)\), the actual through
interval is covered by the negative-saddle endpoint, a compact family
of strict source passages, and the high-height interval. Consecutive
pieces overlap in actual trajectories. The completed transition has
negative curvature on the endpoint and compact middle, and bounded
first derivative on the high-height interval. All cuts and estimates
are uniform on the fixed \(p\)-interval.
\end{lemma}

\subsection{The equatorial roots and the strong foliation}
\label{ann:subsec-negative-saddle}

With \(a=-1\), write \eqref{ann:source-field} as
\[
 X'=-(M+LY),\qquad Y'=-Y(XB+\zeta Y),
\]
where
\begin{equation}\label{ann:source-coefficients}
\begin{aligned}
 M&=1-pX+3\gamma X^2+sX^3,& B&=1+2\gamma+sX,\\
 L&=1-A(\gamma+sX)+C(\gamma+sX)^2,&
 \zeta&=Cs(\gamma+sX).
\end{aligned}
\end{equation}

\begin{lemma}\label{ann:lem-axis-roots}
The roots of \(M\) are simple and have the order
\(-R<x_+<x_f\), where \(R\to\infty\), \(x_+\to p^{-1}\).
For \(M_-(q)=M(-q)\), \(B_-(q)=B(-q)\), put
\(\kappa=-M_-'(R)\) and \(\rho_-=\kappa/(RB_-(R))\).
Then
\begin{equation}\label{ann:negative-root-identity}
 sR=3\gamma+p/R+R^{-2},
\end{equation}
and
\begin{equation}\label{ann:small-negative-index}
 \kappa\ge cp_0,\qquad
 0<\rho_-\le C(R^{-1}+|\gamma|+sR)\longrightarrow0 .
\end{equation}
At \(x_+\), the rates \(L_+=-M'(x_+)\), \(\nu_+=x_+B(x_+)\)
give \(0<\rho_0\le\rho_+=\nu_+/L_+\le\rho_1<\infty\).
\end{lemma}
\begin{proof}
On compact sets \(M\to1-pX\), yielding \(x_+\) and the negative
sign immediately to its right. Since \(M(0)>0\), the cubic has a
second positive root and a negative root, both escaping uniformly.
Substitution at \(-R\) proves \eqref{ann:negative-root-identity}.
The factorization
\[
 M=s(X+R)(X-x_+)(X-x_f),\qquad sRx_+x_f=1
\]
gives
\(\kappa=x_+^{-1}(1+x_+/R)(1+R/x_f)\ge cp_0\).
Direct differentiation also gives
\(\rho_-=(2p/R+3\gamma+3/R^2)/B_-(R)\).
The positive rates converge on the compact \(p\)-interval.
\end{proof}

Fix \(Q\gg P\), then \(R_0>2Q\), and shrink the coefficients so that
\begin{equation}\label{ann:delta-Q}
 d_Q=Q^{-2}+P/Q+|\gamma|+sR\ll1 .
\end{equation}
In \(q=-X\) the negative tail is
\begin{equation}\label{ann:negative-tail-field}
 q'=M_-(q)+L_-(q)Y,\qquad
 Y'=Y\{qB_-(q)+\zeta_-(q)Y\},
\end{equation}
with
\[
 \begin{aligned}
 M_-&=1+pq+3\gamma q^2-sq^3,& B_-&=1+2\gamma-sq,\\
 L_-&=1-A(\gamma-sq)+C(\gamma-sq)^2,&
 \zeta_-&=Cs(sq-\gamma).
 \end{aligned}
\]
Let \(\mathcal V\) denote \eqref{ann:negative-tail-field}.  Use the
positive normal time \(d\tau_n=qB_-(q)\,dt\), and write
\(V=\mathcal V/(qB_-(q))\).  In \(u=\log q\), \(w=Y/q^2\),
put \(c=L_-/B_-\) and \(k=\zeta_-q/B_-\).  The axis speed is
\begin{equation}\label{ann:negative-slow-speed}
 h_-(q)=\frac{M_-(q)}{q^2B_-(q)}
       =\frac{q^{-2}+p/q+3\gamma-sq}{B_-(q)}=O(d_Q).
\end{equation}
The exact normal-time equations are
\[
 u'=h_-+cw,\qquad w'=w\{1-2h_-+(k-2c)w\}.
\]
The functions \(h_-\), \(c-1\), \(k\) and their first five Euler
derivatives \(D_q^j=(q\partial_q)^j\) are \(O(d_Q)\) on
\(Q/2\le q\le2R\) and fixed logarithmic collars.
Extend this triple consistently as in
Lemma~\ref{lem:normal-time-projection}.  Its hypotheses hold with
\(\epsilon=Cd_Q\); choosing \(Q\) first and then the small
coefficient bounds makes this quantity sufficiently small.  The
projection is therefore constructed directly from the negative-tail
field, independently of either return-map estimate.

It gives a strong projection \(\pi(q,Y)\) for \(V\) with uniform
scaled bounds
on the whole corridor, including the root \(R\).  The induction
through one further order, using the fifth coefficient derivatives,
gives a bounded \(C^4\) norm for \(\log\pi-\log q\) in
\((\log q,Y/q^2)\).  Throughout this
section, the footpoint refers to this normal-time projection.
In particular, its invariance equation is
\begin{equation}\label{ann:strong-projection}
 \pi(q,0)=q,\qquad \pi_*V=b_-(r)\partial_r,\qquad
 b_-(r)=\frac{M_-(r)}{rB_-(r)},\qquad r=\pi(q,Y).
\end{equation}

Use the section
\begin{equation}\label{ann:negative-strong-section}
 Y/q^2=w_0,\qquad 0<w_0<1/4 .
\end{equation}
At the limiting corner, the leaf over \(r\) satisfies
\(Y-q^2/2=-r^2/2\), and meets the section transversely at
\(q=r(1-2w_0)^{-1/2}\).  The uniform projection estimates give a
unique intersection over every \(R_0\le r<R\), of the form
\[
 q_s(r)=rU_s(r),\qquad U_s\asymp1,\qquad q_s'(r)\asymp1,
 \qquad |(r\partial_r)^j\log U_s|\le C\quad(1\le j\le3).
\]
We use this footpoint \(r\) to parametrize the strong section.

Figure~\ref{fig:tails-roots} shows the compact core, long tails
and root configuration.
% Schematic domains and equatorial roots for a=-1, p in [p_0,P].
% The arrows in panel (a) denote positive-time transition maps.
% Panel (b) is an ordered root diagram, not an orbit on the equator.
\begin{figure}[htb]
\centering
\setlength{\unitlength}{1pt}
\begin{picture}(410,182)
 \put(205,172){\makebox(0,0){\small (a) The transition on fixed physical sections}}
 \put(45,148){\makebox(0,0){\(\Sigma_-\)}}
 \put(155,148){\makebox(0,0){\(X=-L\)}}
 \put(255,148){\makebox(0,0){\(X=L\)}}
 \put(365,148){\makebox(0,0){\(\Sigma_+\)}}
 \linethickness{.6pt}
 \put(142,128){\vector(-1,0){85}}
 \put(242,128){\vector(-1,0){74}}
 \put(353,128){\vector(-1,0){86}}
 \put(100,113){\makebox(0,0){\small outgoing tail}}
 \put(205,113){\makebox(0,0){\small compact passage}}
 \put(310,113){\makebox(0,0){\small incoming tail}}
 \put(45,97){\makebox(0,0){\(\dfrac{-\delta-\gamma}{s}\)}}
 \put(365,97){\makebox(0,0){\(\dfrac{\delta-\gamma}{s}\)}}
 \put(205,69){\makebox(0,0){\small (b) Equatorial roots in the same coordinate}}
 \linethickness{.35pt}
 \put(24,42){\vector(1,0){368}}
 \put(400,42){\makebox(0,0){\(X\)}}
 \put(195,39){\line(0,1){6}}
 \put(195,29){\makebox(0,0){\small \(0\)}}
 \put(74,42){\circle*{3}}
 \put(219,42){\circle*{3}}
 \put(338,42){\circle*{3}}
 \put(74,54){\makebox(0,0){\(-R\)}}
 \put(219,54){\makebox(0,0){\(x_+\)}}
 \put(338,54){\makebox(0,0){\(x_f\)}}
 \put(74,18){\makebox(0,0){\small \(-R\to-\infty\)}}
 \put(219,18){\makebox(0,0){\small \(x_+\to p^{-1}\)}}
 \put(338,18){\makebox(0,0){\small \(x_f\to+\infty\)}}
 \put(74,0){\makebox(0,0){\small \(\rho_-\to0\)}}
 \put(219,0){\makebox(0,0){\small \(\rho_0\le\rho_+\le\rho_1\)}}
\end{picture}
\caption{Two consequences of the rescaling \(v=\gamma+sX\),
\(z=s^2Y\), for \(\alpha_0=-s^2\) and \(p\in[p_0,P]\).
The positions and distances are schematic. In (a), a fixed compact
core does not include the tails leading to \(\Sigma_\pm\);
the arrows denote the three factors of the positive-time transition.
In (b), the ordered zeros of \(M\) are shown as \(s,\gamma\to0\).
The negative incident saddle escapes and its characteristic number
vanishes, although the positive incident saddle remains on a compact
scale. The axis arrow gives the orientation of \(X\), not the flow.}
\label{fig:tails-roots}
\end{figure}


\subsection{An integrating coordinate up to the negative saddle}

The linear normal equation along the negative axis is
\(d\log Y/dq=qB_-(q)/M_-(q)\). Its integral suggests the coordinates
\begin{equation}\label{ann:negative-action}
 \mathcal A_-(r)=\int_Q^r\frac{qB_-(q)}{M_-(q)}\,dq,\qquad
 \mathcal D_-=\frac{M_-(r)}{rB_-(r)}\partial_r .
\end{equation}
Then \(\mathcal D_-\mathcal A_-=1\), and
\(\mathcal A_-\to\infty\) at \(R\).
Thus \(\mathcal A_-\) covers the entire interval up to the saddle,
although the upper bound \(R\) for the footpoint tends to infinity.
The next lemma relates this linear integral to the height
of the actual orbit on the fixed cut \(q=Q\).

\begin{lemma}\label{ann:lem-global-action-preparation}
The height \(Y_Q(r)\) at the fixed cut \(q=Q\) of the orbit through
the strong section over \(r\) satisfies
\begin{equation}\label{ann:negative-action-preparation}
 \log Y_Q=c_\theta-\mathcal A_-(r)+\mathcal R_-(r),
 \qquad R_0\le r<R,
\end{equation}
where, for \(j=1,2\),
\begin{equation}\label{ann:negative-action-remainder}
 |\mathcal D_-^j\mathcal R_-|\le C(d_Q+e_{R_0}),\qquad
 e_{R_0}=C_QR_0^2e^{-c(R_0-Q)/P}\longrightarrow0.
\end{equation}
The value of \(\mathcal R_-\) need not be small.
\end{lemma}
\begin{proof}
Apply Lemma~\ref{lem:normal-time-action-height} to the triple just
constructed. The polynomial in \eqref{ann:negative-tail-field}
gives \(M_-(q)/q\ge p_0/2\) on \([Q/2,Q]\) after shrinking the
coefficients, and \(b_->0\) up to its simple root \(R\).
After \(R_0\) is fixed, the same polynomial gives
\(b_-\le CP\) on \([Q,R_0]\). Thus the lemma constructs the
whole actual height interval within \(Q/2<q<2R\), and yields
\begin{equation}\label{ann:global-action-exact}
 \log Y_Q=-\mathcal A_-(r)+2\log r+\log\mathcal U(r)+E_Q(Y_Q),
\end{equation}
with the unit and fixed-cut bounds stated there. In particular,
\begin{equation}\label{ann:negative-small-height}
 0<Y_Q(r)\le e_{R_0}=C_QR_0^2e^{-c(R_0-Q)/P}.
\end{equation}
Choose \(R_0\) large and then shrink the coefficient neighborhood.
Taking \(c_\theta=0\) and
\(\mathcal R_-=2\log r+\log\mathcal U+E_Q(Y_Q)\),
\eqref{eq:normal-time-action-jets} gives
\(\mathcal D_-\mathcal R_-=1-K\) and
\(\mathcal D_-^2\mathcal R_-=-\mathcal D_-K\), proving the
claimed bounds. The lemma also gives \(Y_Q\downarrow0\) at \(R\).
\end{proof}

\subsection{Derivatives on the return sections}

We now continue from the strong section to the actual return section,
retaining the footpoint power and the divided closing. In the following
proof, \(w_0\) keeps the small-lip restriction
\(0<w_0<1/4\) of \eqref{ann:negative-strong-section}; only
\(w_1\) is chosen in \((1/2,1)\).

\begin{lemma}\label{ann:lem-negative-gate}
Let \(\Gamma_-(r)\) denote the resulting coordinate on the negative
return section.
Then
\begin{equation}\label{ann:negative-gate-derivative}
 -\Gamma_-'(r)=\Omega_-(\theta)r^{1-\lambda_\gamma}U_-(r)>0,
 \qquad \Omega_->0,
\end{equation}
where \(U_-\) and its reciprocal are bounded and
\begin{equation}\label{ann:negative-gate-jets}
 |\mathcal D_-^j\log U_-|\le C(d_Q+e_{R_0}),\qquad j=1,2 .
\end{equation}
\end{lemma}
\begin{proof}
Fix \(w_0\), \(w_1\in(1/2,1)\), \(\mu=3/2\), and the small
physical width \(\delta\) before its nested output windows. All phase
derivatives fix coefficients and multiplier. On the strong section the
inverse of \(\upsilon=\widehat\pi(u,w_0)\), \(\upsilon=\log r\), is
\(u=a(\upsilon)=\log q_s(r)\). The projection bounds give
\begin{equation}\label{ann:repair-inverse-start}
 a_1\asymp1,\qquad |a_2|+|a_3|\le C,\qquad 1<q_s/r<2,
 \qquad a_j=\partial_\upsilon^j a.
\end{equation}
Also, for \(h=M_-/(r^2B_-)\),
\begin{equation}\label{ann:repair-action-budget}
 \mathcal D_-=h\partial_\upsilon,\qquad
 |h|+|h_\upsilon|\le Cd_Q.
\end{equation}
The scalar graph and its logarithmic-radius equation are
\begin{equation}\label{ann:negative-scalar-graph}
 Y_q=f=\frac{Y(qB_-+\zeta_-Y)}{M_-+L_-Y},\qquad
 W_u=F(u,W)=W\left\{\frac{B_-+ZW}{D+L_-W}-2\right\},
\end{equation}
where \(W=Y/q^2\), \(D=q^{-2}+p/q+3\gamma-sq\), and \(Z=q\zeta_-\).
On \(w_0\le W\le w_1\), after choosing \(Q\) and shrinking the
coefficients,
\begin{equation}\label{ann:repair-outgoing-strip}
 D+L_-W\ge c,\qquad F(u,w_0)>c,\qquad F(u,w_1)<-c,\qquad F_W\le-\mu.
\end{equation}
The needed coefficient and state derivatives are bounded; explicitly,
\begin{equation}\label{ann:repair-rational-identities}
 qf_Y=F_W+2,\qquad
 F_{WW}=\frac{2D(ZD-B_-L_-)}{(D+L_-W)^3},\qquad
 F_{WWW}=-\frac{6DL_-(ZD-B_-L_-)}{(D+L_-W)^4}.
\end{equation}
Thus the strip is invariant and \(q'>0\). Since
\(\sup_{r<R}q_s(r)<q_H=(\delta+\gamma)/s\), every orbit reaches
this raw exit once, even when \(q_s>R\).

To extract the long-tail power, compare with
\[
 F^0_\gamma(W)=W\left\{\frac{1+2\gamma}{3\gamma+L_\gamma W}-2\right\},
 \quad L_\gamma=1-A\gamma+C\gamma^2,
 \quad w_\gamma=\frac{1-4\gamma}{2L_\gamma}.
\]
Here \(F^0_\gamma(w_\gamma)=0\),
\((F^0_\gamma)'(w_\gamma)=-2+\lambda_\gamma\),
\((F^0_\gamma)''=O(|\gamma|)\), and
\(F-F^0_\gamma=O(sq+p/q+q^{-2})\) in \(C^1_W\).
Scalar contraction, with \(q_0=q_s(r)\), consequently gives
\begin{equation}\label{ann:repair-memory-remainder}
 \left|f_Y-\frac{\lambda_\gamma}{q}\right|
 \le C\left(s+pq^{-2}+q^{-3}
       +\frac{|\gamma|}{q}(q_0/q)^\mu\right).
\end{equation}
The last term retains entrance memory and has integral at most
\(|\gamma|/\mu\), independently of the tail length. Hence
\begin{equation}\label{ann:repair-integrated-remainder}
 I(r)=\int_{q_0}^{q_H}(f_Y-\lambda_\gamma/q)\,dq,\qquad
 |I(r)|\le C(\delta+|\gamma|+p/q_0+q_0^{-2}).
\end{equation}
For the fixed-entry first variation,
\begin{equation}\label{ann:negative-first-variation}
 (J_Y)_q=f_YJ_Y,
\end{equation}
so the exact cocycle is
\begin{equation}\label{ann:negative-free-factor}
 \frac{J_Y(q)}{J_Y(q_0)}=(q/q_0)^{\lambda_\gamma}
 \exp\int_{q_0}^{q}\{f_Y(\widetilde q,Y(\widetilde q))-\lambda_\gamma/\widetilde q\}\,d\widetilde q.
\end{equation}

The moving initial value \(W(a(\upsilon);\upsilon)=w_0\) gives
\begin{equation}\label{ann:repair-moving-variation}
 J:=\partial_\upsilon W=-F(a,w_0)a_1
 \exp\int_a^uF_W\,dv<0,\qquad |J|\le Ce^{-\mu(u-a)}.
\end{equation}
Differentiate its logarithm using the triangular variation argument of
Lemma~\ref{lem:scalar-normalization}. The first integral forcing is
\(F_{WW}J\); in the second row it is
\(F_{WWW}J^2+F_{WW}J\partial_\upsilon\log(-J)\).
The moving lower-limit terms involve only \(a_1,a_2,a_3\) and the
bounded entrance derivatives of \(F\). Integrating the decaying
forcings in order yields
\begin{equation}\label{ann:repair-moving-log-first}
 |\partial_\upsilon\log(-J)|\le C,
\end{equation}
\begin{equation}\label{ann:repair-moving-log-second}
 |\partial_\upsilon^2\log(-J)|\le C.
\end{equation}
No reciprocal of \(J\) is estimated separately.

For raw admission choose \(0<\eta\ll\delta\), then require
\(2sR<\eta\). Every orbit crosses \(sq=\eta\) in the strip.
Contraction on the fixed remaining raw interval, compared with
\(w_\Gamma(x)=(\sqrt{1+x^2}+1)^{-1}\), gives
\[
 |W_H-w_\Gamma(\delta+\gamma)|
 \le C(\eta/(\delta+\gamma))^\mu
 +C_{\eta,\delta}(|A-1|+|C-1|+|\gamma|+sp+s^2).
\]
The raw height is \(z_H=(\delta+\gamma)^2W_H\); its organizing
height shifts by \(O(\delta|\gamma|)\). Choose \(\eta/\delta\)
from the fixed window reserve, then shrink the coefficients. On this
whole admitted image the divided adapter \(\chi=\mathcal C_{\theta,m}\)
has the ordinary \(C^3\) bounds of Lemma~\ref{phys:lem-c3-adapters}.
Set \(\Gamma_-(r)=\chi(z_H(r))\). Its observer identities and the
preceding two rows give
\begin{equation}\label{ann:repair-completed-log-jets}
 |\partial_\upsilon\log(-\Gamma_-')|
 +|\partial_\upsilon^2\log(-\Gamma_-')|\le C.
\end{equation}
Finally separate \(F_W=-2+\lambda_\gamma+(qf_Y-\lambda_\gamma)\)
in the moving cocycle:
\begin{equation}\label{ann:negative-raw-gate-factor}
 -z_H'=\Omega_\theta r^{1-\lambda_\gamma}\widehat U_-,\qquad
 \Omega_\theta=s^{2-\lambda_\gamma}(\delta+\gamma)^{\lambda_\gamma},
\end{equation}
\begin{equation}\label{ann:repair-unit-factor}
 \widehat U_-=F(a,w_0)a_1(q_s/r)^{2-\lambda_\gamma}e^{I(r)},\qquad
 U_-=\chi'(z_H)\widehat U_-.
\end{equation}
The factors \(\widehat U_-,U_-\) are positive units, with
\(\Omega_\theta\) kept separate. Subtracting
\((1-\lambda_\gamma)\upsilon+\log\Omega_\theta\) in
\eqref{ann:repair-completed-log-jets} bounds two ordinary derivatives
of \(\log U_-\). Equation~\eqref{ann:repair-action-budget} then gives
\(\mathcal D_-\log U_-=O(d_Q)\) and
\(\mathcal D_-^2\log U_-=O(d_Q^2)\), proving the assertion.
The same bounds before the adapter give
\(\mathcal D_-^jz_H=O((sr)^{2-\lambda_\gamma}d_Q^j)\), \(j=1,2\).
\end{proof}

At the other saddle the characteristic number remains compact.
Its Dulac estimate must still include the actual incoming coordinate,
because the curvature is taken in the physical affine input.

\begin{lemma}[The fixed-cut height and the actual positive half-gate]
\label{ann:p0-positive-interface}
For fixed \(p\in[p_0,P]\Subset(0,\infty)\), physical windows,
and then fixed \(Q\), there is a common positive radius \(H_*\)
such that the map from the local positive-saddle exit to \(q=Q\)
is \(\Psi_\theta(v)\), with \(\Psi_\theta(0)=0\), bounded
\(C^3\) norm, and derivative bounded above and away from zero.
The actual incoming half-gate \(x=\Gamma_+(H)\), on a slightly
larger interval than \([0,H_*]\), satisfies
\[
 \Gamma_+'>0,\qquad
 \left|\frac{\Gamma_+''}{\Gamma_+'}\right|
 +\left|\frac{\Gamma_+'''}{\Gamma_+'}\right|\le C_Q.
\]
Its exact value-and-slope normalization and its inverse on a smaller
normalized image have uniform \(C^3\) bounds. Its full raw image
is admitted by the preassigned incoming window.
For every fixed \(0<\overline H\le H_*\),
\(y_+(\overline H)=Y_Q^+(\overline H)\) has a uniform positive
lower bound. No uniform lower bound for the raw \(\Gamma_+'\)
or bound for the raw inverse is asserted.
\end{lemma}
\begin{proof}
The fixed rescaled axis segment from the outgoing saddle cut to \(q=Q\)
has a positive speed margin. Its finite hitting map \(\Psi\) therefore
has \(\Psi(0)=0\), positive derivative units and uniform \(C^3\)
bounds; no raw-tail scale enters this map.

For the incoming half-map, follow the stable separatrix backwards. In
the limiting field its denominator \(D=Y+1-pX\) cannot first vanish,
since \(D_X=XY/D-p>0\) near a proposed zero; a quadratic upper
barrier excludes finite blow-up. The compact half-family consequently
enters a common parabolic band at a fixed \(X=L_0\). A buffered chord
has a finite hitting map \(A_\theta(H)=Y(L_0)/L_0^2\) with positive
derivative units and uniform \(C^3\) bounds, including \(H=0\).

On the outward tail \(F_w\le-\lambda<0\), and the state derivatives
through order three are bounded. The relative-jet argument
\eqref{tail:relative-jet-equations}, with
\(0<J_1\le e^{-\lambda t}\), bounds \(J_2/J_1,J_3/J_1\).
Exact value-and-slope normalization thus gives the asserted map and
inverse bounds. The fixed-port argument of
Lemma~\ref{ann:p0-fixed-tail-ports} admits its whole raw chord: that
argument needs only this common parabolic band, retains the entrance
memory and \(O(d|\gamma|)\) port shift, and reaches the original
window through its fixed raw tube. Compose the ordinary bridge on this
admitted image.

Finally Lemma~\ref{lem:uniform-positive-dulac} and \(\Psi\) give
\(y_+(H)=c_+H^{\rho_+}e^{\mathcal R_+(H)}\), with \(c_+\) a
positive unit. Thus for fixed \(0<\overline H<1\),
\(y_+(\overline H)\ge c_Q\overline H^{\rho_1}e^{-M_Q}>0\).
This height bound involves no lower bound on the raw incoming slope.
\end{proof}

\begin{lemma}\label{ann:lem-positive-dulac}
There is a fixed incoming coordinate \(H>0\) at \(x_+\) for which
the height on \(q=Q\) is
\begin{equation}\label{ann:positive-dulac-map}
 Y_Q^+(H)=c_+H^{\rho_+}e^{\mathcal R_+(H)}.
\end{equation}
For \(0<\epsilon<\min(1,\rho_0)\), the Euler derivatives of
\(\mathcal R_+\) through order two are bounded by
\(CH^\epsilon(1+|\log H|)\).
The continuation to the input \(x\) in \eqref{eq:return}
has the form \(x=\Gamma_+(H)\), with
\begin{equation}\label{ann:positive-gate-derivative}
 \Gamma_+'(H)=\Omega_+(\theta)U_+(H)>0,
\end{equation}
where \(U_+\) and its reciprocal are bounded and its first two
Euler logarithmic derivatives have the same vanishing bound.
\end{lemma}
\begin{proof}
Apply Lemma~\ref{lem:uniform-positive-dulac} to the rates in
Lemma~\ref{ann:lem-axis-roots}, then compose with the fixed-cut map
\(\Psi\) above. Hadamard division gives
\(\Psi(v)=\Psi'(0)ve^{j(v)}\), \(j(0)=0\), with bounded \(C^2\)
norm. Its first two Euler rows along the Dulac map are
\(O(H^{\rho_+})\), so this composition preserves the asserted remainder.
For the incoming half-gate set
\(\Omega_+=\Gamma_+'(0)\), \(U_+=\Gamma_+'(H)/\Gamma_+'(0)\).
The preceding relative \(C^3\) bounds give
\(\mathcal E\log U_+,\mathcal E^2\log U_+=O(H)\),
\(\mathcal E=H\partial_H\). The final input differs from raw height
only by the translation \(a_\theta\).
\end{proof}

\subsection{Endpoint curvature and the total count}

\begin{lemma}[Exhaustion of the low central heights]
\label{ann:p0-low-height-exhaustion}
There are fixed \(R_a<R_b\), \(R_a>2Q\), and \(u_L>0\), chosen
before the final coefficient shrink, such that every actual through
orbit with \(0<u:=Y(0)\le2u_L\) is represented uniquely by the
same negative strong footpoint \(R_b<r<R\) and by the positive
Dulac coordinate \(0<H<\overline H\). The identification on the
fixed cut is
\[
 N_\theta(u)=y_\theta(r)=Y_Q^+(H),
 \qquad N_\theta(u):=Y(Q;Y(0)=u).
\]
All actual input and output windows are those already fixed; zero
height is not included as a through orbit.
\end{lemma}
\begin{proof}
Start at \((Q,y)\), \(0<y<w_0Q^2\). Until the strong-section hit,
the footpoint \(p_f=\pi(q,Y)\) increases below \(R\), and the
projection estimates give
\[
 (\log q)'\ge h(\log p_f)+(1-C\epsilon)w>0,\qquad
 q<2p_f<2R,\qquad (\log w)'\ge\nu_0>0.
\]
Thus there is one transverse hit in finite normal time, inside the
unchanged corridor; it may have \(q_s>R\). Scalar first variation
and \(\partial_q(Y-w_0q^2)=q_sF(\log q_s,w_0)>0\) give
\(dq_s/dy<0\), hence \(dr/dy<0\). The reverse construction in
Lemma~\ref{lem:normal-time-action-height} has
\(y_r=-Ky/b<0\), \(K\ge1/2\), and \(y\to0\) at \(R\).
It is therefore a bijection \([R_a,R)\to(0,y(R_a)]\), identical
to the forward construction by uniqueness.

On \(0\le q\le Q\), choose a small signed height strip with
\(M_-+L_-Y>c>0\). Its finite map satisfies
\(N(0)=0\), \(c_Q\le N'\le C_Q\), with common finite-order bounds.
Choose \(\overline H\) for the positive Euler errors, then \(R_a\)
so that \(e_{R_a}<\inf y_+(\overline H)\) and the fixed-cut errors
are small; take \(R_0=R_a\). On every fixed footpoint interval the
exact action-height formula bounds the unit and action, giving
\(\inf y(R_a)>0\). Choose \(R_b>R_a\) with
\(e_{R_b}<\inf y(R_a)\), then require \(R>2R_b\).
The same fixed-interval formula gives
\[
 y(R_b)=R_b^2\mathcal U(R_b)
 e^{-\mathcal A(R_b)+E_Q(y(R_b))}\ge\eta_b>0.
\]
Choose \(2u_L\) below the domain of \(N\) and \(\eta_b/C_Q\).
Then \(N(u)<y(R_b)\) selects a unique \(r>R_b\); it is also below
\(y_+(\overline H)\) and selects a unique positive Dulac preimage.
Equality at \(q=Q\) identifies the orbit, so both previously admitted
gates and the divided closing apply throughout this interval.
\end{proof}

The exhaustion lemma identifies the same orbit in the two endpoint
coordinates. Equating their descriptions of \(Y_Q\) gives
\begin{equation}\label{ann:endpoint-matching}
 \rho_+\log H+\mathcal R_+(H)
       =-\mathcal A_-(r)+\mathcal R_-(r)+c_\theta .
\end{equation}
Set \(\mathcal D=H\partial_H\) and
\(K_-=1-\mathcal D_-\mathcal R_-\ge1/2\).
Then
\begin{equation}\label{ann:rH}
 Hr_H=-\frac{\rho_++\mathcal D\mathcal R_+}{K_-}
                  \frac{M_-(r)}{rB_-(r)}<0 .
\end{equation}
By \eqref{ann:negative-small-height} and
\eqref{ann:positive-dulac-map}, throughout this endpoint
\[
 H\le C_QR_0^{2/\rho_0}e^{-c(R_0-Q)/(P\rho_1)}.
\]
Thus the same choice of \(R_0\) makes all positive Dulac remainders
small on the entire interval \(R_0\le r<R\).

\begin{proposition}\label{ann:prop-endpoint-curvature}
The complete map
\(x=\Gamma_+(H)\mapsto\Pi_\theta(x)=\Gamma_-(r(H))\)
has \(\Pi_\theta''(x)<0\) on this endpoint. Its displacement has
at most two zeros counted with multiplicity.
\end{proposition}
\begin{proof}
Write \(y(H)=Y_Q^+(H)\) and
\(\mathcal O(y(r))=\Gamma_-(r)\). Put
\(b=M_-/(rB_-)\), \(\mathcal D_-=b\partial_r\), and
\(K=-\mathcal D_-\log y\ge1/2\). The action-height jets give \(K=1+O(d_Q+e_{R_a})\) and
\(\mathcal D_-K=O(d_Q+e_{R_a})\). The coefficient quotients are
\begin{align}
 M_-'/(rB_-)&=(p/r+6\gamma-3sr)/B_-=O(d_Q),
 \label{ann:Mprime-quotient}\\
 M_-/(r^2B_-)&=(r^{-2}+p/r+3\gamma-sr)/B_-=O(d_Q),
 \label{ann:M-quotient}\\
 sM_-/(rB_-^2)&=O(d_Q).\label{ann:B-quotient}
\end{align}
Thus \(b'=O(d_Q)\). The identity
\(\mathcal O'=(-\Gamma_-')b/(yK)\), together with
Lemma~\ref{ann:lem-negative-gate}, gives
\[
 e_{\mathcal O}:=1+y\frac{\mathcal O''}{\mathcal O'}
 =-\frac{\mathcal D_-\log(-\Gamma_-')+b'-\mathcal D_-\log K}{K},
 \qquad |e_{\mathcal O}|\le Cd_Q+C_Qe_{R_a}.
\]
No reciprocal negative saddle rate occurs.
Set \(\mathsf A=\rho_++H\mathcal R_+'\). In the actual input
\(x=\Gamma_+(H)\), the full chain rule is
\begin{equation}\label{ann:complete-log-curvature}
 H\partial_H\log\Pi_x'
 =-1+\mathsf A e_{\mathcal O}
      +H\partial_H\log\mathsf A-H\partial_H\log\Gamma_+'.
\end{equation}
Choose \(Q\) first. For this fixed cut choose \(\overline H\)
so that \(|\mathsf A-\rho_+|\le\min(1,\rho_0/2)\) and the last
two terms each have absolute value at most \(1/16\).
Choose \(R_a\) using the positive fixed-cut image margin, as in
Lemma~\ref{ann:p0-low-height-exhaustion}, and then shrink the
coefficients so that
\( |e_{\mathcal O}|\le1/[16(\rho_1+1)]\).
The right side is at most \(-13/16\). Since
\(\Pi_x'>0\), \(H\Gamma_+'>0\),
\[
 \frac{\Pi_{xx}}{\Pi_x'}
 =\frac{H\partial_H\log\Pi_x'}{H\Gamma_+'}<0.
\]
For the physical multiplier \(m>0\), \(\Delta''=-m\Pi''>0\);
Rolle's theorem with multiplicities gives the endpoint count.
\end{proof}

\begin{lemma}[Fixed parabolic tails and the physical sections]
\label{ann:p0-fixed-tail-ports}
Fix the physical sections \(\Sigma_\pm=\{v=\pm\delta\}\), their
nested windows, and \(0<d<\delta\le1/32\), with \(d\le1/100\).
On an enlarged compact strict source family, one can choose a fixed
\(L\), and then sufficiently small bounds on the raw coefficients
and on \(s>0\), so that the two outward graphs reach
\(v=\gamma\pm d\). Their entire images are admitted by uniformly
regular increasing maps
\[
 \mathcal B_{\sigma,\theta}:z(\gamma+\sigma d)
       \longmapsto z(\sigma\delta),\qquad \sigma=\pm1.
\]
These maps and their inverses have uniform phase \(C^3\) bounds on
slightly larger fixed windows. The divided closing map is defined on
the full negative image for \(0\le m\le m_{\rm ad}\), where
\(m_{\rm ad}>0\) is fixed before any zero-counting multiplier bound.
\end{lemma}
\begin{proof}
Choose \(|\gamma|<\min\{d/2,(\delta-d)/2,1/500\}\).
For \(r=Le^t\), \(\xi=sr\), write \(w=Y/r^2\) on either tail.
Its exact equation is
\[
 \dot w=F_\sigma(t,w)
 =w\left(\frac{n_{0,\sigma}+n_{1,\sigma}w}
                   {\ell_\sigma w+d_\sigma}-2\right),
\]
where
\[
\begin{aligned}
 v_\sigma&=\gamma+\sigma\xi,
 &\ell_\sigma&=1-Av_\sigma+Cv_\sigma^2,\\
 d_\sigma&=3\gamma+\sigma\xi-\sigma p/r-a/r^2,
 &n_{0,\sigma}&=1+2\gamma+\sigma\xi,
 &n_{1,\sigma}&=\sigma C\xi v_\sigma.
\end{aligned}
\]
The entrance \(L\) is chosen on the enlarged limiting family first.
Continuous dependence on the fixed core then puts both entrance
values strictly inside \([1/4,3/4]\). Increase \(L\), if necessary,
so that \(|p|/L\le1/40\) and \(|a|/L^2\le1/4000\), and require
\(A,C\in[9/10,11/10]\) and \(sL<d\). The fixed core is retained
when these choices are made. On
\[
 1/4\le w\le3/4,\quad \xi\le1/100,
 \quad |\gamma|\le1/500,
 \quad |p|/r\le1/40,
 \quad |a|/r^2\le1/4000
\]
the denominator is bounded below and the vector field points strictly
inward at both edges. Differentiating this rational quotient, as in
\eqref{ann:repair-rational-identities}, gives
\begin{equation}\label{ann:p0-tail-envelope}
 F_w\le-\lambda,\qquad \lambda=7/5,
 \qquad |F_{ww}|\le10,
 \qquad |F_{ww}|\le240|d_\sigma|.
\end{equation}
Thus the band remains invariant up to
\(T=\log(d/(sL))\), independently of its length.

The invariant band controls the graph throughout the long tail. To
compose with the regular transitions, we must also place its whole image
inside the fixed physical windows.
At the organizing coefficients the exact radial graph is
\[
 w_\Gamma(\xi)=\frac1{\sqrt{1+\xi^2}+1},\qquad
 \xi\partial_\xi w_\Gamma=F^0_\sigma(\xi,w_\Gamma).
\]
On the preceding coefficient box, the mean-value formula gives
\[
 |F_\sigma-F^0_\sigma|
 \le100\bigl(|\gamma|+|A-1|\xi+|C-1|\xi^2
                         +|p|/r+|a|/r^2\bigr).
\]
For \(T=\log(d/(sL))\), contraction gives the uniform error
\[
 \begin{split}
 \mathcal E_\theta(L,T)={}&e^{-\lambda T}
       |w_\sigma(0)-w_\Gamma(sL)|\\
 &+100\int_0^T e^{-\lambda(T-t)}
 \left(|\gamma|+|A-1|sLe^t+|C-1|(sLe^t)^2
       +\frac{|p|}{L}e^{-t}+\frac{|a|}{L^2}e^{-2t}\right)dt.
 \end{split}
\]
In particular \(|w_\sigma(T)-w_\Gamma(d)|\le\mathcal E_\theta\).
This retains the entrance memory; \(1<\lambda<2\) bounds both
decaying-power convolutions. For fixed \(L\) and bounded \(a,p\),
\(\mathcal E_\theta\to0\) uniformly as
\(s,\gamma,A-1,C-1\to0\).
With \(g(v)=\sqrt{1+v^2}-1\), the actual port heights satisfy
\[
 |z_\sigma-g(\gamma+\sigma d)|
       \le d^2\mathcal E_\theta(L,T)+2d|\gamma|.
\]

On the fixed annulus \(d/2\le|v|\le2\delta\), the organizing graph
has \(-D_0(v,g(v))\ge v^2/4\ge d^2/16\).
Choose a fixed tube about this graph, and then shrink the raw
coefficients so that \(-D_\theta\ge d^2/32\) there. The actual
graph equation
\[
 z_v=\frac{z(R_\theta(v)+Cvz)}
                {[1-Av+Cv^2]z-P_\theta(v)}
\]
has bounded state derivatives through order three in the tube.
Gronwall's inequality, with the port error just obtained, gives
admission to that tube by a first-exit argument. Its finite physical
length then gives the asserted regular maps and inverse bounds.
First place their entire images strictly inside the prescribed inner
windows. Next choose \(m_{\rm ad}\) so that
\(m d_\theta(K_-^\sharp)\) lies strictly inside the fixed domain of
\(f_\theta\). Only after these domain assertions do we compose the
divided closing map.
\end{proof}

\begin{lemma}[Actual curvature on a compact positive source block]
\label{ann:p0-middle-actual-curvature}
Fix \(0<p_0\le p\le P\) and a buffered positive central-height
interval \(0<u_-\le u\le u_+\). For \(a=-1\), after a common
coefficient shrink the complete transition has \(\Pi_{xx}<0\) on
this whole block in the actual affine input. The assertion is
uniform for \(0\le m\le m_{\rm ad}\).
\end{lemma}
\begin{proof}
We first construct the limiting strict block. On the two half-graphs
\(X=\sigma r\),
\[
 Y_r=\frac{rY}{Y+1-\sigma pr},\qquad Y(0)=u.
\]
The height increases and is bounded above by
\(u_++1+r^2+2Pr\). The negative denominator is automatically
positive. On the positive half, at a sufficiently small positive
denominator its derivative \(rY/D-p\) is positive, which excludes
a first zero. More quantitatively,
\[
 D\ge\min\{1/2,u_-/(4P^2)\}>0.
\]
These inequalities exclude finite blow-up and give both infinite
half-graphs. A common parabolic entrance is fixed as follows. For
\(Q=\max\{4,16P\}\), \(w=Y/r^2\), and \(t=\log r\), the source
equation satisfies
\[
 0<w\le1/4\ \Longrightarrow\ (\log w)_t\ge2/3,
 \qquad w\ge3/4\ \Longrightarrow\ (\log w)_t\le-2/5.
\]
Its entrance bounds at \(Q\) are uniform. After a fixed additional
logarithmic time both half-graphs enter the strict parabolic band.

The resulting enlarged strict family is admitted on the actual
sections by Lemma~\ref{ann:p0-fixed-tail-ports}. On this compact
family, Proposition~\ref{ann:prop-source-curvature} gives
\(S''=-p\mathcal I_{bb}\) with a uniform negative margin since
\(p\ge p_0\). The positive affine normalization of \(S\) has the
same margin up to bounded positive factors.
Proposition~\ref{prop:source-normalization} with phase order two
therefore gives \(\Phi''<0\) after a common coefficient shrink.

Retain the exact centres and gains of the actual maps \(h,k\) in
\eqref{eq:reference-normalization}, including both bridges and divided
closing. With \(c_h=h'(b_0)>0\), \(c_k=k'(b_0)>0\), one has
\begin{equation}\label{ann:p0-two-actual-adapters}
 \Pi_{xx}(x)=\frac{c_k}{c_h^2}\Phi''\!\left(
 b_0+\frac{x+a_\theta-h(b_0)}{c_h}\right)<0.
\end{equation}
The curvature sign is preserved by the positive prefactor, even when
both absolute gains tend to zero.
\end{proof}

\begin{lemma}[Paired first variations at high central height]
\label{ann:p0-paired-high-derivative}
For the fixed compact positive-ratio region, there are fixed
\(K_H,z_*,C_H>0\) such that the complete actual transition has
\(0<\Pi_x'\le C_H\) on
\[
 u=Y(0)>K_H(1+p^2),\qquad 0<s^2u<z_*.
\]
The constant is independent of the tested phase point. The actual
multiplier restriction \(mC_H<1/2\) gives \(\Delta'>1/2\).
\end{lemma}
\begin{proof}
The derivative assertion is Proposition~\ref{sharp:upper-high} with
\(a=-1\). Its paired absolute variations
\eqref{sharp:upper-tail-ratio} cancel the common
\(s^{2-\lambda_\gamma}\) power before the actual bridges and
ordinary divided closing are composed. For the high range freeze
\(H=s^2u\); the signed normalization in
Corollary~\ref{rew:cor-high-derivative-signed} gives the same
\(C_H\) at every phase point.

For the domain assertion, apply
Lemma~\ref{ann:p0-fixed-tail-ports} to that signed compact source
block, fixing its core, tail entrance, ports and raw-bridge tubes
before shrinking the coefficients and \(z_*\). This admits both
full raw images, including the positive-height equality and
double-zero faces. Shrink \(z_*\) further so that
\[
 \sigma_\theta\le-D_\theta(\gamma,H)
   =H(\ell_\gamma-\bar a)\le2H<d_{\rm out}.
\]
Restrict the receiver to \(H<\min\{h_{\rm out},z_*\}\), and
choose \(m<\min\{m_{\rm ad},1/(2C_H)\}\) last.
\end{proof}

\begin{proof}[Proof of Lemma~\ref{ann:lem-compact-positive-domains}]
Fix the coefficient interval, finite smoothness, physical windows and
closing-domain cap \(m_{\rm ad}\). Lemma~\ref{ann:p0-low-height-exhaustion}
chooses \(Q,\overline H,R_a=R_0,R_b,u_L\) in that order and matches
all \(0<u\le2u_L\) through the fixed-cut equality
\(N(u)=y(r)=Y_Q^+(H)\). Its full raw images are admitted.
Independently fix the core, tail entrance and raw tubes for
Lemma~\ref{ann:p0-paired-high-derivative}, obtaining \(K_H,C_H,z_*\).
Choose \(u_H>\max\{2u_L,K_H(1+P^2)\}\) and apply
Lemma~\ref{ann:p0-middle-actual-curvature} on \([u_L/2,3u_H]\).
Intersect these coefficient bounds, including \(R>2R_b\), without
changing \(Q\).

Choose the unscaled cap and coefficient bounds so that
\[
 0<h_{\rm cap}\le z_*,\qquad
 L_{\rm raw}h_{\rm cap}<d_{\rm out}/4,\qquad
 |\alpha_0|<d_{\rm out}/4,
 \quad 0<\ell_\theta(\gamma)\le L_{\rm raw}.
\]
Then every through orbit with \(H=z(\gamma)<h_{\rm cap}\) satisfies
\(\sigma_\theta\le-D_\theta(\gamma,H)
=-\alpha_0+\ell_\theta(\gamma)H<d_{\rm out}/2\).
Let \(J\) be the original buffered receiver clipped by
\(H<\min\{h_{\rm cap},h_{\rm out}\}\), in the affine coordinate
\(x=z_{\rm in}-a_\theta\). The final global choice takes
\(h_{\rm out}\le h_{\rm cap}\), as in Section~\ref{sec:conclusion}.

The strictly increasing central height
\(u(x)=z_\theta(\gamma;a_\theta+x)/s^2\) gives the relative open
intervals
\[
 I_E=J\cap\{0<u<2u_L\},\qquad
 I_M=J\cap\{u_L<u<2u_H\},\qquad
 I_H=J\cap\{u>u_H\}.
\]
They cover \(J\), including all four internal equality levels, and
have physical overlaps \(u_L<u<2u_L\) and \(u_H<u<2u_H\).
The endpoint estimate holds on the larger \(r>R_a\) domain, the
middle estimate on \([u_L/2,3u_H]\), and the high estimate below
\(z_*\). Original window and speed clips preserve this cover; the
fixed-cut match has its own image margin. Raw overlap lengths need no
positive lower bound.

Uniqueness identifies these as restrictions of the same completed map
\eqref{eq:return}. Proposition~\ref{ann:prop-endpoint-curvature} and
Lemma~\ref{ann:p0-middle-actual-curvature} give \(\Pi_{xx}<0\) on
\(I_E\cup I_M\), while
Lemma~\ref{ann:p0-paired-high-derivative} gives
\(0<\Pi_x'\le C_H\) on \(I_H\). All cuts, caps and estimates
precede the final multiplier restriction and geometric capture.
\end{proof}
